% ========== DROP-IN REPLACEMENT FOR SECTION 3 OF project_proposal.tex ==========
% Replace the entire \section{The Solution} ... \subsection{Deliverables} ... \end{itemize}
% with the contents below (from \section{The Solution} through the final \end{itemize}).

\section{The Solution}

\subsection{Framing: Forecasting + Anomaly Detection + Decision Support}
We deliver three components to maximize operational impact and demonstrate advanced Python at scale: (1)~\textbf{Forecasting} --- next-hour meter-level energy consumption using out-of-core, leakage-safe features and a baseline vs.\ LightGBM comparison. (2)~\textbf{Anomaly detection} --- residual-based scoring (actual minus predicted) to flag malfunctioning meters or abnormal consumption, implemented efficiently over large arrays. (3)~\textbf{Decision support} --- a ranked output (e.g., top meters or buildings to audit) so operators know where to act first.

\subsection{Data Processing Strategy (50M+ Rows)}
\textbf{Out-of-core by default.} The main train table is never fully loaded into memory. We use chunked reads (e.g., 1--2M rows per chunk) with pandas or Dask; merge each chunk with in-memory building metadata and weather; then either write features to disk (partitioned by site or time) or feed a downstream sampler. Peak memory target: under 4\,GB per process. In-memory: building metadata, weather (per site), and any per-site or per-meter subset used for a single model.

\subsection{Feature Engineering (Time-Aware, Leakage-Safe)}
Features use only past and current information. \textit{Time:} hour, day of week, month, is\_weekend, holiday (calendar). \textit{Lags:} same hour yesterday, same day last week. \textit{Rolling:} 24h and 168h rolling mean with an expanding or correctly aligned window so no future data leaks. \textit{Building:} square footage, primary use, year built. \textit{Weather:} air temperature (and optionally lagged). Implementation: vectorized NumPy/pandas per chunk; no Python loops over rows. Unit tests assert that no feature depends on future timestamps.

\subsection{Modeling Strategy}
\textbf{Baseline:} Per-meter mean or ``same hour yesterday'' predictor; time-based train/validation split (e.g., last 3--6 months held out). \textbf{Stronger model:} LightGBM (handles large data, categoricals, fast training). One global model with site/meter identifiers as categoricals, or per-site models trained in parallel. Target: continuous meter consumption; metric: RMSE or CV-RMSE. Success target: LightGBM achieves at least 20--40\% RMSE improvement over baseline on the same split.

\subsection{Scalability Strategy}
\textbf{Partitioning:} by site (19 sites) or by time chunks. \textbf{Parallelization:} multiprocessing or concurrent.futures --- e.g., one worker per site for feature build or training, or parallel chunks for feature construction. No shared mutable state; each worker reads its partition and writes outputs. Benchmark: report runtime and speedup (e.g., 1 vs.\ 4 workers) on a fixed subset (e.g., 5M rows).

\subsection{Acceleration Strategy}
\textbf{Numba JIT} on the hottest path: e.g., anomaly scoring over millions of residuals (score per meter or per timestep) or a custom rolling aggregation not easily expressed in pandas. Use \texttt{@jit(nopython=True)}; compare runtime vs.\ pure Python loop. If a bottleneck is in pandas/NumPy, vectorization is preferred before JIT. No Cython unless profiling shows a clear need.

\subsection{Profiling and Benchmarking}
\textbf{Tools:} \texttt{cProfile} for CPU hotspots, \texttt{memory\_profiler} for peak memory. \textbf{Metrics:} runtime (s), peak memory (GB), speedup (optimized/baseline). \textbf{Comparisons:} (1)~data load: naive full read (or OOM) vs.\ chunked; (2)~feature build: single-thread vs.\ vectorized and/or parallel; (3)~anomaly: Python loop vs.\ Numba or vectorized. One benchmark table in the report with realistic numbers; at least one bottleneck identified and improved.

\subsection{Engineering Strategy}
\textbf{Structure:} \texttt{src/} (load, features, model, anomaly, decision\_support), \texttt{scripts/}, \texttt{config/}, \texttt{tests/}. \textbf{CLI:} one entry point, e.g.\ \texttt{python -m src.cli run --config config.yaml}. \textbf{Config:} YAML/TOML for paths, \texttt{chunk\_size}, \texttt{seed}, model params. \textbf{Logging:} standard library or structlog; INFO level. \textbf{Typing:} type hints on public functions. \textbf{Tests:} pytest; at least feature-leakage check and one small end-to-end run. \textbf{Reproducibility:} \texttt{requirements.txt}, fixed seed in config, README with run instructions.

\subsection{Benchmark Table (Targets)}
\begin{table}[h]
\centering
\small
\begin{tabular}{@{}lccccc@{}}
\toprule
\textbf{Component} & \textbf{Baseline} & \textbf{Optimized} & \textbf{Runtime} & \textbf{Peak mem.} & \textbf{Success target} \\
\midrule
Data load & Full read / OOM & Chunked & --- & $<$4\,GB & No OOM \\
Feature build (5M rows) & Single-thread & Vectorized $\pm$ parallel & --- & --- & $\geq 2\times$ speedup \\
Forecasting & Mean predictor & LightGBM & --- & --- & 20--40\% RMSE gain \\
Anomaly (1M rows) & Python loop & Numba / vectorized & --- & --- & $\geq 5\times$ speedup \\
\bottomrule
\end{tabular}
\caption{Benchmark template: fill with measured runtimes and memory.}
\end{table}

\subsection{Deliverables}
\begin{itemize}[noitemsep]
    \item Reproducible pipeline (scripts + config) producing forecasts, anomaly scores, and a decision-support list.
    \item Benchmark table with baseline vs.\ optimized runtimes and memory; at least one documented speedup.
    \item Final report ($\leq$4 pages): problem, data, solution (three components), benchmarks, results, limitations.
    \item Code: typed, tested (leakage check + one integration test), CLI, README.
\end{itemize}
